\documentclass[11pt,a4paper]{article}
% Wspólny szablon dla kart supermocy (jedna karta = jedna strona A4)
% Użycie w karcie: \documentclass[11pt,a4paper]{article} % Wspólny szablon dla kart supermocy (jedna karta = jedna strona A4)
% Użycie w karcie: \documentclass[11pt,a4paper]{article} % Wspólny szablon dla kart supermocy (jedna karta = jedna strona A4)
% Użycie w karcie: \documentclass[11pt,a4paper]{article} \input{../preambula} \begin{document}

% Silnik tectonic (XeTeX): polskie znaki przez fontspec, nie przez inputenc/T1
\usepackage{fontspec}
\setmainfont{Arial}
\usepackage[polish]{babel}
\usepackage[a4paper,margin=1.5cm]{geometry}
\usepackage{amsmath,amssymb}
\usepackage{tikz}
\usetikzlibrary{arrows.meta,positioning,calc,shapes.geometric,decorations.pathreplacing}
\usepackage[most]{tcolorbox}
\usepackage{enumitem}
\usepackage{booktabs}
\usepackage{xcolor}

\pagestyle{empty}
\setlength{\parindent}{0pt}
\setlength{\parskip}{0.4em}

\definecolor{cel}{HTML}{1F4E79}
\definecolor{regula}{HTML}{E8F1FA}
\definecolor{pulapka}{HTML}{FDECEA}
\definecolor{okres}{HTML}{C0392B}
\definecolor{skonczony}{HTML}{1E8449}

% Ramka na regułę
\newtcolorbox{regula}[1][]{
  colback=regula, colframe=cel, boxrule=0.8pt, arc=2pt,
  left=6pt, right=6pt, top=4pt, bottom=4pt, title=#1,
  fonttitle=\bfseries, coltitle=white
}

% Ramka na pułapkę z kartkówki
\newtcolorbox{pulapka}[1][]{
  colback=pulapka, colframe=okres, boxrule=0.8pt, arc=2pt,
  left=6pt, right=6pt, top=4pt, bottom=4pt, title=#1,
  fonttitle=\bfseries, coltitle=white
}

% Nagłówek karty: numer, tytuł, cel
\newcommand{\naglowekKarty}[3]{%
  {\LARGE\bfseries\color{cel} Karta #1: #2}\\[2pt]
  {\large Cel: #3}\par
  \vspace{0.3em}\noindent\rule{\linewidth}{1pt}\par
}

% Pole na obliczenia ucznia: ramka z delikatną kratką, argument = wysokość
\newcommand{\miejsce}[1]{%
  \par\vspace{0.2em}\noindent
  \begin{tikzpicture}
    \draw[step=0.5cm, gray!25, very thin] (0,0) grid (\textwidth,#1);
    \draw[gray!60, rounded corners=2pt] (0,0) rectangle (\textwidth,#1);
    \node[anchor=north west, font=\scriptsize, text=gray] at (0.05,#1) {obliczenia};
  \end{tikzpicture}\par}
 \begin{document}

% Silnik tectonic (XeTeX): polskie znaki przez fontspec, nie przez inputenc/T1
\usepackage{fontspec}
\setmainfont{Arial}
\usepackage[polish]{babel}
\usepackage[a4paper,margin=1.5cm]{geometry}
\usepackage{amsmath,amssymb}
\usepackage{tikz}
\usetikzlibrary{arrows.meta,positioning,calc,shapes.geometric,decorations.pathreplacing}
\usepackage[most]{tcolorbox}
\usepackage{enumitem}
\usepackage{booktabs}
\usepackage{xcolor}

\pagestyle{empty}
\setlength{\parindent}{0pt}
\setlength{\parskip}{0.4em}

\definecolor{cel}{HTML}{1F4E79}
\definecolor{regula}{HTML}{E8F1FA}
\definecolor{pulapka}{HTML}{FDECEA}
\definecolor{okres}{HTML}{C0392B}
\definecolor{skonczony}{HTML}{1E8449}

% Ramka na regułę
\newtcolorbox{regula}[1][]{
  colback=regula, colframe=cel, boxrule=0.8pt, arc=2pt,
  left=6pt, right=6pt, top=4pt, bottom=4pt, title=#1,
  fonttitle=\bfseries, coltitle=white
}

% Ramka na pułapkę z kartkówki
\newtcolorbox{pulapka}[1][]{
  colback=pulapka, colframe=okres, boxrule=0.8pt, arc=2pt,
  left=6pt, right=6pt, top=4pt, bottom=4pt, title=#1,
  fonttitle=\bfseries, coltitle=white
}

% Nagłówek karty: numer, tytuł, cel
\newcommand{\naglowekKarty}[3]{%
  {\LARGE\bfseries\color{cel} Karta #1: #2}\\[2pt]
  {\large Cel: #3}\par
  \vspace{0.3em}\noindent\rule{\linewidth}{1pt}\par
}

% Pole na obliczenia ucznia: ramka z delikatną kratką, argument = wysokość
\newcommand{\miejsce}[1]{%
  \par\vspace{0.2em}\noindent
  \begin{tikzpicture}
    \draw[step=0.5cm, gray!25, very thin] (0,0) grid (\textwidth,#1);
    \draw[gray!60, rounded corners=2pt] (0,0) rectangle (\textwidth,#1);
    \node[anchor=north west, font=\scriptsize, text=gray] at (0.05,#1) {obliczenia};
  \end{tikzpicture}\par}
 \begin{document}

% Silnik tectonic (XeTeX): polskie znaki przez fontspec, nie przez inputenc/T1
\usepackage{fontspec}
\setmainfont{Arial}
\usepackage[polish]{babel}
\usepackage[a4paper,margin=1.5cm]{geometry}
\usepackage{amsmath,amssymb}
\usepackage{tikz}
\usetikzlibrary{arrows.meta,positioning,calc,shapes.geometric,decorations.pathreplacing}
\usepackage[most]{tcolorbox}
\usepackage{enumitem}
\usepackage{booktabs}
\usepackage{xcolor}

\pagestyle{empty}
\setlength{\parindent}{0pt}
\setlength{\parskip}{0.4em}

\definecolor{cel}{HTML}{1F4E79}
\definecolor{regula}{HTML}{E8F1FA}
\definecolor{pulapka}{HTML}{FDECEA}
\definecolor{okres}{HTML}{C0392B}
\definecolor{skonczony}{HTML}{1E8449}

% Ramka na regułę
\newtcolorbox{regula}[1][]{
  colback=regula, colframe=cel, boxrule=0.8pt, arc=2pt,
  left=6pt, right=6pt, top=4pt, bottom=4pt, title=#1,
  fonttitle=\bfseries, coltitle=white
}

% Ramka na pułapkę z kartkówki
\newtcolorbox{pulapka}[1][]{
  colback=pulapka, colframe=okres, boxrule=0.8pt, arc=2pt,
  left=6pt, right=6pt, top=4pt, bottom=4pt, title=#1,
  fonttitle=\bfseries, coltitle=white
}

% Nagłówek karty: numer, tytuł, cel
\newcommand{\naglowekKarty}[3]{%
  {\LARGE\bfseries\color{cel} Karta #1: #2}\\[2pt]
  {\large Cel: #3}\par
  \vspace{0.3em}\noindent\rule{\linewidth}{1pt}\par
}

% Pole na obliczenia ucznia: ramka z delikatną kratką, argument = wysokość
\newcommand{\miejsce}[1]{%
  \par\vspace{0.2em}\noindent
  \begin{tikzpicture}
    \draw[step=0.5cm, gray!25, very thin] (0,0) grid (\textwidth,#1);
    \draw[gray!60, rounded corners=2pt] (0,0) rectangle (\textwidth,#1);
    \node[anchor=north west, font=\scriptsize, text=gray] at (0.05,#1) {obliczenia};
  \end{tikzpicture}\par}


\begin{document}

\naglowekKarty{2}{Jaki to procent i ile wynosi całość}{zamienić „część z całości” na procent i odwrotnie}

\vspace{0.3em}
\begin{regula}[REGUŁA – DWA PYTANIA I DWA RUCHY]
\begin{enumerate}[leftmargin=*, itemsep=3pt]
  \item \textbf{Jaki to procent?} Zapisz ułamek $\dfrac{\text{część}}{\text{całość}}$ i rozszerz go do mianownika $100$. Licznik to szukany procent. \quad $\dfrac{5}{25}=\dfrac{20}{100}=20\%$.
  \item \textbf{Ile wynosi całość?} Znajdź $1\%$ (podziel część przez liczbę procentów), potem pomnóż przez $100$. \quad $20\%=14\ \Rightarrow\ 1\%=0{,}7\ \Rightarrow\ 100\%=70$.
\end{enumerate}
\end{regula}

\vspace{0.4em}
\textbf{Krok po kroku: 5 uczniów z 25 ma zielone oczy. Jaki to procent?}
\begin{enumerate}[leftmargin=*, itemsep=3pt]
  \item Część to $5$, całość to $25$: $\dfrac{5}{25}$.
  \item Rozszerz do mianownika $100$ (mnożymy przez $4$): $\dfrac{5\cdot 4}{25\cdot 4}=\dfrac{20}{100}=20\%$.
\end{enumerate}
\textbf{Druga droga:} podziel i pomnóż przez $100$. $5:25=0{,}2$, a $0{,}2\cdot 100\%=20\%$. Gdy mianownik nie jest dzielnikiem $100$, ta droga jest szybsza: $\dfrac{3}{8}=0{,}375=37{,}5\%$.

\begin{center}
\begin{tikzpicture}[>=Stealth]
  % Siatka 5x5: 25 uczniów, 5 zaznaczonych
  \foreach \r in {0,...,4} {
    \foreach \c in {0,...,4} {
      \draw (\c*0.8,-\r*0.8) rectangle ++(0.8,-0.8);
    }
  }
  \foreach \c in {0,...,4} { \fill[cel!45] (\c*0.8,0) rectangle ++(0.8,-0.8); \draw (\c*0.8,0) rectangle ++(0.8,-0.8); }
  \node[font=\small\bfseries, color=cel, anchor=west] at (4.4,-0.4) {$5$ zielonych oczu};
  \node[font=\small, anchor=west] at (4.4,-1.4) {całość: $25$ uczniów};
  \node[font=\small\bfseries, color=okres, anchor=west] at (4.4,-2.4) {$\dfrac{5}{25}=\dfrac{20}{100}=20\%$};
  \node[font=\small, anchor=west] at (4.4,-3.3) {jeden rząd z pięciu to $\dfrac{1}{5}=20\%$};
\end{tikzpicture}
\end{center}

\vspace{-0.2em}
\textbf{Przykłady rozstrzygnięte:}
\begin{center}
\renewcommand{\arraystretch}{1.5}
\begin{tabular}{@{}lll@{}}
\toprule
Zadanie & Obliczenie & Wynik \\
\midrule
$6$ z $20$ & $\dfrac{6}{20}=\dfrac{30}{100}$ & $30\%$ \\[4pt]
$9$ z $12$ & $\dfrac{9}{12}=\dfrac{3}{4}$ & $75\%$ \\[4pt]
$3$ z $8$ & $3:8=0{,}375$ & $37{,}5\%$ \\
$10\%$ całości to $30\ \text{zł}$ & $30\cdot 10$ & $300\ \text{zł}$ \\
$25\%$ całości to $12$ & $12\cdot 4$ (bo $25\%=\dfrac{1}{4}$) & $48$ \\
$20\%$ całości to $14$ & $1\%=0{,}7$, \ $0{,}7\cdot 100$ & $70$ \\
\bottomrule
\end{tabular}
\end{center}

\begin{pulapka}[PUŁAPKI Z KARTKÓWKI]
\begin{enumerate}[leftmargin=*, itemsep=2pt]
  \item Czytaj, o co pyta zadanie. Z $20$ kostek zjedzono $6$: zjedzono $30\%$, ale \emph{zostało} $14$ kostek, czyli $70\%$.
  \item Całość stoi w mianowniku. Ile procent liczby $50$ to $20$? $\dfrac{20}{50}=40\%$, a nie $\dfrac{50}{20}=250\%$.
  \item Jeśli $10\%$ to $30$, całość to $300$ (razy $10$), a nie $3$.
\end{enumerate}
\end{pulapka}

\end{document}
